\section{Kapitel 5 -- Implementierung}\label{kapitel-5-implementierung}

\begin{center}\rule{0.5\linewidth}{0.5pt}\end{center}

\subsection{5.1 Überblick und
Implementierungsstrategie}\label{uxfcberblick-und-implementierungsstrategie}

Dieses Kapitel dokumentiert die konkrete Umsetzung des in Kapitel 4
konzipierten Systems. Im Fokus stehen die technischen Entscheidungen auf
Code-Ebene: die interne Struktur des FastAPI-Backends mit seinen
Schichten aus Modellen, Repositories, Services und Routern; die
KI-Generierungspipeline von der Ereigniserkennung bis zum
Ticker-Eintrag; die Frontend-Implementierung mit React,
Context-basiertem Zustandsmanagement und dem slash-command-getriebenen
Eingabesystem; sowie die n8n-Workflow-Implementierung der fünf zentralen
Automatisierungsworkflows. Ergänzend werden die TypeScript-Migration und
die Testabdeckung als Qualitätsmerkmale des Entwicklungsprozesses
dokumentiert.

Die Implementierung folgt dem Grundsatz \textbf{``so einfach wie nötig,
so modular wie sinnvoll''}: Standardpfade (CRUD, Listabfragen) sind
bewusst direkt gehalten (Router → Repository), während die
LLM-Generierung mit ihren Retry-, Semaphore- und Few-Shot-Mechanismen in
einen dedizierten Service extrahiert ist.

Die n8n-Workflows (Abschnitt 5.5) werden abweichend von der
konzeptionellen Reihenfolge in Kapitel 4 nach der
Frontend-Implementierung beschrieben: Die Workflows bauen auf
Backend-Endpunkten auf, die bis dahin vollständig dokumentiert sind,
sodass die Beschreibung der Orchestrierungslogik erst an dieser Stelle
vollständig nachvollziehbar ist.

\begin{center}\rule{0.5\linewidth}{0.5pt}\end{center}

\subsection{5.2 Backend-Implementierung}\label{backend-implementierung}

\subsubsection{5.2.1 Projektstruktur und Application Entry
Point}\label{projektstruktur-und-application-entry-point}

Das Backend ist unterhalb von \passthrough{\lstinline!backend/app/!} in
sieben funktionale Pakete gegliedert:

\begin{lstlisting}
backend/app/
├── api/v1/          # HTTP-Router (FastAPI)
├── models/          # SQLAlchemy ORM-Modelle (17 Dateien)
├── repositories/    # Datenbankzugriff (je Entität eine Klasse)
├── schemas/         # Pydantic-Schemas (Request / Response)
├── services/        # Fachliche Services (llm_service, ticker_service)
├── utils/           # Hilfsmodule (llm_context_builders, evaluation_metrics)
└── core/            # Konfiguration, Datenbankverbindung, Enums, Konstanten
\end{lstlisting}

Diese Aufteilung entspricht dem klassischen
Repository-Service-Router-Muster: Router delegieren an Repositories für
einfache CRUD-Operationen und an Services für Logik, die mehrere
Repositories oder externe Aufrufe koordiniert.

\passthrough{\lstinline!app/main.py!} ist der einzige Einstiegspunkt. Er
registriert alle 12 Router unter dem Pfad-Prefix
\passthrough{\lstinline!/api/v1!}, konfiguriert CORS-Middleware und
bindet statische Dateien (Thumbnail-Cache) ein. Der WebSocket-Endpunkt
\passthrough{\lstinline!/ws/media!} wird ohne API-Prefix eingehängt, da
er kein REST-Ressourcenpfad ist. Beim Start führt
\passthrough{\lstinline!Base.metadata.create\_all(bind=engine)!} ein
idempotentes Schema-Bootstrapping durch, das bei erstmaliger
Inbetriebnahme alle Tabellen anlegt; Schema-Änderungen im
Produktionsbetrieb werden via Alembic-Migrationen gesteuert.

Die Anwendung stellt zwei Metaendpunkte bereit:

{\def\LTcaptype{none} % do not increment counter
\begin{longtable}[]{@{}ll@{}}
\toprule\noalign{}
Endpunkt & Funktion \\
\midrule\noalign{}
\endhead
\bottomrule\noalign{}
\endlastfoot
\passthrough{\lstinline!GET /!} & Statusantwort mit Version
(\passthrough{\lstinline!0.3.0!}) \\
\passthrough{\lstinline!GET /health!} & Datenbankverbindungscheck
(\passthrough{\lstinline!healthy!}/\passthrough{\lstinline!degraded!}) \\
\end{longtable}
}

\begin{center}\rule{0.5\linewidth}{0.5pt}\end{center}

\subsubsection{5.2.2 Datenmodelle: Datenbankschema und
API-Validierung}\label{datenmodelle-datenbankschema-und-api-validierung}

\textbf{SQLAlchemy-Datenbankmodelle} --- Das Schema umfasst \textbf{17
ORM-Modelle} (18 Datenbanktabellen, davon eine via Migration). Die
zentralen Domänenobjekte sind:

{\def\LTcaptype{none} % do not increment counter
\begin{longtable}[]{@{}
  >{\raggedright\arraybackslash}p{(\linewidth - 4\tabcolsep) * \real{0.1848}}
  >{\raggedright\arraybackslash}p{(\linewidth - 4\tabcolsep) * \real{0.2065}}
  >{\raggedright\arraybackslash}p{(\linewidth - 4\tabcolsep) * \real{0.6087}}@{}}
\toprule\noalign{}
\begin{minipage}[b]{\linewidth}\raggedright
Modell
\end{minipage} & \begin{minipage}[b]{\linewidth}\raggedright
Tabelle
\end{minipage} & \begin{minipage}[b]{\linewidth}\raggedright
Kernfelder
\end{minipage} \\
\midrule\noalign{}
\endhead
\bottomrule\noalign{}
\endlastfoot
\passthrough{\lstinline!Country!} & \passthrough{\lstinline!countries!}
& \passthrough{\lstinline!id!}, \passthrough{\lstinline!name!},
\passthrough{\lstinline!uid!} \\
\passthrough{\lstinline!Team!} & \passthrough{\lstinline!teams!} &
\passthrough{\lstinline!id!}, \passthrough{\lstinline!uid!},
\passthrough{\lstinline!name!}, \passthrough{\lstinline!external\_id!},
\passthrough{\lstinline!is\_partner\_team!} \\
\passthrough{\lstinline!Competition!} &
\passthrough{\lstinline!competitions!} & \passthrough{\lstinline!id!},
\passthrough{\lstinline!title!}, \passthrough{\lstinline!season\_id!} \\
\passthrough{\lstinline!Match!} & \passthrough{\lstinline!matches!} &
\passthrough{\lstinline!id!}, \passthrough{\lstinline!uid!},
\passthrough{\lstinline!match\_state!},
\passthrough{\lstinline!match\_phase!},
\passthrough{\lstinline!ticker\_mode!} \\
\passthrough{\lstinline!Event!} & \passthrough{\lstinline!events!} &
\passthrough{\lstinline!id!}, \passthrough{\lstinline!match\_id!},
\passthrough{\lstinline!event\_type!}, \passthrough{\lstinline!time!},
\passthrough{\lstinline!description!} \\
\passthrough{\lstinline!SyntheticEvent!} &
\passthrough{\lstinline!synthetic\_events!} &
\passthrough{\lstinline!id!}, \passthrough{\lstinline!match\_id!},
\passthrough{\lstinline!type!}, \passthrough{\lstinline!data!}
(JSONB) \\
\passthrough{\lstinline!TickerEntry!} &
\passthrough{\lstinline!ticker\_entries!} & → siehe unten \\
\passthrough{\lstinline!Lineup!} & \passthrough{\lstinline!lineups!} &
\passthrough{\lstinline!player\_id!}, \passthrough{\lstinline!status!},
\passthrough{\lstinline!formation\_place!},
\passthrough{\lstinline!position!} \\
\passthrough{\lstinline!MatchStatistic!} &
\passthrough{\lstinline!match\_statistics!} & 19 typisierte Spalten
(Pässe, Zweikämpfe, Schüsse\ldots) \\
\passthrough{\lstinline!PlayerStatistic!} &
\passthrough{\lstinline!player\_statistics!} &
\passthrough{\lstinline!player\_id!}, \passthrough{\lstinline!rating!},
\passthrough{\lstinline!goals!}, \passthrough{\lstinline!assists!},
\ldots{} \\
\passthrough{\lstinline!StyleReference!} &
\passthrough{\lstinline!style\_references!} &
\passthrough{\lstinline!event\_type!},
\passthrough{\lstinline!instance!}, \passthrough{\lstinline!league!},
\passthrough{\lstinline!text!} (Few-Shot) \\
\passthrough{\lstinline!MediaQueue!} &
\passthrough{\lstinline!media\_queue!} &
\passthrough{\lstinline!media\_id!},
\passthrough{\lstinline!thumbnail\_url!},
\passthrough{\lstinline!event\_id!}, \passthrough{\lstinline!status!} \\
\passthrough{\lstinline!MediaClip!} &
\passthrough{\lstinline!media\_clips!} &
\passthrough{\lstinline!source!}, \passthrough{\lstinline!vid!},
\passthrough{\lstinline!thumbnail\_url!},
\passthrough{\lstinline!match\_id!} \\
\end{longtable}
}

Das \passthrough{\lstinline!TickerEntry!}-Modell hat besondere Bedeutung
im Lifecycle des Systems: Es verknüpft über
\passthrough{\lstinline!event\_id!} (FK auf
\passthrough{\lstinline!events!}) und
\passthrough{\lstinline!synthetic\_event\_id!} (FK auf
\passthrough{\lstinline!synthetic\_events!}) den generierten Text mit
seinem Auslöser. Die Felder \passthrough{\lstinline!status!}
(\passthrough{\lstinline!draft!}/\passthrough{\lstinline!published!}/\passthrough{\lstinline!rejected!})
und \passthrough{\lstinline!source!}
(\passthrough{\lstinline!ai!}/\passthrough{\lstinline!manual!}) bilden
den Ticker-Lifecycle ab (vgl. Kap. 4.3.2). Drei zusammengesetzte Indizes
(\passthrough{\lstinline!ix\_ticker\_match\_status!},
\passthrough{\lstinline!ix\_ticker\_match\_phase!},
\passthrough{\lstinline!ix\_ticker\_event\_id!}) optimieren die
häufigsten Abfragen: publizierte Einträge je Spiel, Phase-Filter und den
Deduplizierungs-Check per Event. Das Status-Enum wird als nicht-nativ
gespeichert (\passthrough{\lstinline!native\_enum=False!}), damit
Migrationen ohne Enum-Typ-Änderungen in PostgreSQL auskommen.

\textbf{Pydantic-Schemas und API-Validierung} --- Für jeden Router gibt
es eigenständige Pydantic-Schemas für Create-, Update- und
Response-Strukturen. Alle Schemas nutzen
\passthrough{\lstinline!alias\_generator=to\_camel!} aus
\passthrough{\lstinline!pydantic.alias\_generators!}, damit
JSON-Payloads in camelCase kommuniziert werden, während Python-intern
snake\_case gilt.

Ein spezifisches Beispiel illustriert die Alias-Mechanik für
Match-Scores:

\begin{lstlisting}[language=Python]
class MatchResponse(BaseModel):
    home_score: Optional[int] = Field(None, serialization_alias='teamHomeScore')
    away_score: Optional[int] = Field(None, serialization_alias='teamAwayScore')
\end{lstlisting}

Hier wird \passthrough{\lstinline!serialization\_alias!} (nicht
\passthrough{\lstinline!alias!}) verwendet, da das Frontend die
Partner-API-Feldnamen erwartet
(\passthrough{\lstinline!teamHomeScore!}/\passthrough{\lstinline!teamAwayScore!}),
während Schreiboperationen über
\passthrough{\lstinline!populate\_by\_name=True!} auch den Python-Namen
akzeptieren.

Diese duale Alias-Konvention zieht sich durch alle Schemas: Felder, die
Spielstandwerte aus der Partner-API widerspiegeln, erhalten explizite
\passthrough{\lstinline!serialization\_alias!}-Werte mit
Partner-Konvention (\passthrough{\lstinline!teamHomeScore!}).
Standard-Beziehungsfelder hingegen erhalten durch den globalen
\passthrough{\lstinline!alias\_generator=to\_camel!} automatisch
kanonisches camelCase (\passthrough{\lstinline!homeTeamId!},
\passthrough{\lstinline!awayTeamId!}). Die Koexistenz beider
Konventionen in derselben Schema-Klasse ist eine bewusste
Designentscheidung: Das Frontend konsumiert sowohl Partner-API-Daten als
auch eigenständig modellierte Beziehungen, und beide müssen ohne
Feldumbenennungen im Frontend konsistent adressierbar sein.

Für den Match-Status-Lifecycle definiert
\passthrough{\lstinline!MatchUpdate!} ein explizites Alias:

\begin{lstlisting}[language=Python]
match_state: Optional[MatchState] = Field(None, alias='state')
\end{lstlisting}

Für Ticker-Einträge wird der Status über ein separates PATCH-Schema
gesteuert, um versehentliche Massenänderungen durch offene
Update-Schemas zu vermeiden.

\begin{center}\rule{0.5\linewidth}{0.5pt}\end{center}

\subsubsection{5.2.3 Repository-Schicht}\label{repository-schicht}

Jede Entität hat ein dediziertes Repository, das alle SQL-Operationen
kapselt. Alle Repositories erben von einer gemeinsamen
\passthrough{\lstinline!BaseRepository[T]!}-Basisklasse, die die
generische Operation \passthrough{\lstinline!exists()!} bereitstellt;
jede Instanz erhält eine \passthrough{\lstinline!Session!} im
Konstruktor:

\begin{lstlisting}[language=Python]
class TickerEntryRepository(BaseRepository[TickerEntry]):
    model = TickerEntry

    def get_by_event(self, event_id: int) -> Optional[TickerEntry]:
        return self.db.query(TickerEntry).filter(
            TickerEntry.event_id == event_id
        ).first()

    def get_by_match(self, match_id: int, status=None) -> list[TickerEntry]:
        ...
\end{lstlisting}

Diese Kapselung hat zwei Vorteile: Router bleiben frei von ORM-Details,
und Repository-Methoden können direkt mit echten Datenbanken in
Integrationstests geprüft werden --- ohne Mock-Schichten, die
Implementierungsabweichungen verbergen.

Das \passthrough{\lstinline!MatchRepository!} definiert eine interne
\passthrough{\lstinline!\_base\_query()!}-Methode, die Heimteam,
Auswärtsteam, Wettbewerb und Saison per
\passthrough{\lstinline!joinedload!} in einer einzigen Datenbankabfrage
vorlädt. \passthrough{\lstinline!get\_by\_id()!} nutzt diese Methode
konsequent:

\begin{lstlisting}[language=Python]
class MatchRepository(BaseRepository[Match]):
    model = Match

    def _base_query(self) -> Query:
        return self.db.query(Match).options(
            joinedload(Match.home_team),
            joinedload(Match.away_team),
            joinedload(Match.competition),
            joinedload(Match.season),
        )

    def get_by_id(self, match_id: int) -> Optional[Match]:
        return self._base_query().filter(Match.id == match_id).first()
\end{lstlisting}

Diese Methode wird von der KI-Generierungsroute intensiv genutzt, da für
jeden Aufruf vollständige Kontext-Daten (Teamnamen, Wettbewerb) benötigt
werden.

\begin{center}\rule{0.5\linewidth}{0.5pt}\end{center}

\subsubsection{5.2.4 API-Router:
Endpunktstruktur}\label{api-router-endpunktstruktur}

Der Ticker-Bereich ist in drei Router aufgeteilt, die gemeinsam unter
\passthrough{\lstinline!/api/v1/ticker!} eingehängt werden:

{\def\LTcaptype{none} % do not increment counter
\begin{longtable}[]{@{}
  >{\raggedright\arraybackslash}p{(\linewidth - 4\tabcolsep) * \real{0.1753}}
  >{\raggedright\arraybackslash}p{(\linewidth - 4\tabcolsep) * \real{0.2062}}
  >{\raggedright\arraybackslash}p{(\linewidth - 4\tabcolsep) * \real{0.6186}}@{}}
\toprule\noalign{}
\begin{minipage}[b]{\linewidth}\raggedright
Router
\end{minipage} & \begin{minipage}[b]{\linewidth}\raggedright
Datei
\end{minipage} & \begin{minipage}[b]{\linewidth}\raggedright
Kernoperationen
\end{minipage} \\
\midrule\noalign{}
\endhead
\bottomrule\noalign{}
\endlastfoot
\passthrough{\lstinline!ticker\_crud!} &
\passthrough{\lstinline!ticker\_crud.py!} & GET list, GET single, PATCH
status, DELETE, POST manual \\
\passthrough{\lstinline!ticker\_generate!} &
\passthrough{\lstinline!ticker\_generate.py!} & POST
generate/\{event\_id\}, generate-synthetic, generate-batch \\
\passthrough{\lstinline!ticker\_batch!} &
\passthrough{\lstinline!ticker\_batch.py!} & POST translate-batch,
generate-match-phases \\
\end{longtable}
}

Diese Aufteilung verhindert, dass eine einzelne Routerdatei durch die
Kombination aus CRUD-, Generierungs- und Batch-Endpunkten
unübersichtlich wird.

Ein kapselndes Hilfsmuster zieht sich durch alle Router:
\passthrough{\lstinline!require\_or\_404(value, message)!} --- eine
schmale Utility-Funktion in
\passthrough{\lstinline!utils/http\_errors.py!}, die
\passthrough{\lstinline!None!}-Rückgaben aus Repository-Aufrufen direkt
in eine \passthrough{\lstinline!HTTPException(404)!} überführt. Das
eliminiert das repetitive
\passthrough{\lstinline!if obj is None: raise HTTPException(404, ...)!}
in Routern und stellt sicher, dass alle 404-Antworten einheitlich
formuliert sind. Der Batch-Endpunkt
\passthrough{\lstinline!POST /api/v1/ticker/translate-batch/\{match\_id\}!}
ruft für alle publizierten KI-Einträge eines Spiels den LLM-Service mit
reduzierter Temperatur (\passthrough{\lstinline!0.1!}) auf, um
sprachlich deterministische Übersetzungen zu erzeugen --- eine gegenüber
der Originalgenerierung bewusst andere Parametrisierung.

Die wichtigsten Endpunkte gliedern sich in vier funktionale Bereiche:
\textbf{Stammdaten} (Countries, Teams, Matches, Seasons, Competitions),
\textbf{Ticker-Lifecycle} (CRUD, Generierung per Event,
Batch-Generierung für Synthetic Events, Übersetzung), \textbf{Medien \&
Clips} (ScorePlay-Import, Media-Queue, WebSocket) und
\textbf{Spielkontext} (Lineups, Statistiken, Events). Insgesamt umfasst
die API über 70 Routen; die vollständige Endpunktliste ist über die
automatisch generierte OpenAPI-Dokumentation unter
\passthrough{\lstinline!/api/docs!} einsehbar.

\begin{center}\rule{0.5\linewidth}{0.5pt}\end{center}

\subsubsection{5.2.5 Ticker-Service:
Domänenlogik}\label{ticker-service-domuxe4nenlogik}

Der \passthrough{\lstinline!ticker\_service!} kapselt die Domänenlogik
für KI-Einträge. Er koordiniert Datenbankabfragen, Kontextaufbau und den
LLM-Aufruf, hält dabei aber die Schichten klar getrennt:
\passthrough{\lstinline!llm\_service.py!} hat keine
Datenbankabhängigkeit.

Die vier zentralen Funktionen:

\textbf{\passthrough{\lstinline!score\_at\_event()!}} berechnet den
Spielstand zum Zeitpunkt eines Torereignisses. Da die Football-API den
kumulativen Stand liefert, muss der Stand \emph{vor} dem aktuellen Event
rekonstruiert werden. Die Funktion lädt alle Tore bis zur
\passthrough{\lstinline!position!} des Events, unterscheidet zwischen
regulären Toren und Eigentoren (letztere zählen für den Gegner) und
akkumuliert die Scores für Heim- und Auswärtsteam.

\textbf{\passthrough{\lstinline!build\_match\_context()!}} erstellt das
Kontext-Dictionary für den LLM-Prompt: Teamnamen, aktueller Stand,
Matchzustand, Spielminute und Liga.

\textbf{\passthrough{\lstinline!call\_llm()!}} ist der zentrale,
Semaphore-gesicherte LLM-Aufruf. Er lädt zunächst aus dem
\passthrough{\lstinline!StyleReferenceRepository!} bis zu drei
Stilbeispiele für
\passthrough{\lstinline!event\_type + instance + league!} als
Few-Shot-Kontext und delegiert dann an
\passthrough{\lstinline!generate\_ticker\_text()!}:

\begin{lstlisting}[language=Python]
async with _llm_semaphore:   # asyncio.Semaphore(settings.LLM_CONCURRENCY)
    return await generate_ticker_text(
        event_type=event_type,
        style_references=style_references,  # aus DB geladen
        ...
    )
\end{lstlisting}

Die Semaphore ist auf Modulebene als Singleton definiert
(\passthrough{\lstinline!\_llm\_semaphore = asyncio.Semaphore(settings.LLM\_CONCURRENCY)!},
Standard: 8) und begrenzt gleichzeitige LLM-Aufrufe pro Prozessinstanz.

\textbf{\passthrough{\lstinline!make\_ai\_entry()!}} ist ein schlanker
Builder, der aus den LLM-Ergebnissen ein
\passthrough{\lstinline!TickerEntryCreate!}-Schema erzeugt. Er setzt
\passthrough{\lstinline!source='ai'!} und legt den initialen Status
abhängig vom Aufrufkontext fest (\passthrough{\lstinline!draft!} im
Co-op-Modus, \passthrough{\lstinline!published!} im Auto-Modus).

\begin{center}\rule{0.5\linewidth}{0.5pt}\end{center}

\subsubsection{5.2.6 LLM-Service: Prompt-Aufbau und
Provider-Dispatch}\label{llm-service-prompt-aufbau-und-provider-dispatch}

Der \passthrough{\lstinline!LLMService!} ist eine
single-class-Implementierung, die beim Initialisieren den passenden
Client instanziiert und über ein Dispatch-Dictionary aufruft:

\begin{lstlisting}[language=Python]
dispatch = {
    'mock':       self._generate_mock_text,
    'gemini':     self._generate_gemini_text,
    'openrouter': self._generate_openrouter_text,
    'openai':     self._generate_openai_text,
    'anthropic':  self._generate_anthropic_text,
}
return dispatch[self.provider](**kwargs)
\end{lstlisting}

Das Dispatch-Pattern macht es trivial, neue Provider hinzuzufügen:
Implementierung einer Methode
\passthrough{\lstinline!\_generate\_X\_text()!} und ein Eintrag im
Dictionary.

Der Prompt wird in vier modularen Methoden aufgebaut:

\textbf{\passthrough{\lstinline!\_build\_event\_lines()!}} erzeugt die
Fakten-Sektion. Jeder Parameter (Ereignistyp, Detail, Minute, Spieler,
Team) wird nur dann eingebunden, wenn er vorhanden ist. Die
Ereignisbezeichnung wird aus
\passthrough{\lstinline!EVENT\_TYPE\_LABEL!} in lesbares Deutsch
übersetzt (z. B. \passthrough{\lstinline!'goal'!} →
\passthrough{\lstinline!'Tor'!}).

\textbf{\passthrough{\lstinline!\_build\_few\_shot\_block()!}}
formatiert die Stilreferenzen als nummerierte Beispiele:

\begin{lstlisting}
### STILREFERENZEN
Schreibe in exakt diesem Stil (Rhythmus, Wortwahl, Emotionalität):
- 'Mustermann trifft mit einem satten Rechtsschuss – 1:0!'
- 'Eintracht schlägt zu – der erste Treffer des Abends!'
\end{lstlisting}

\textbf{\passthrough{\lstinline!\_build\_prematch\_parts()!}} ergänzt
für Pre-Match-Typen eine zusätzliche harte Instruktion: Das Modell darf
keine Live-Szenen erfinden, sondern ausschließlich Vorschau- und
Analyseinhalte produzieren.

Der vollständig zusammengesetzte System-Prompt folgt der in Abschnitt
4.5.2 beschriebenen modularen Struktur aus sechs Bausteinen
(Rolleninstruktion, Stilinstruktion, Faktenblock, Match-Kontext,
Few-Shot-Block, Regelblock). Für Pre-Match-Typen fügt
\passthrough{\lstinline!\_build\_prematch\_parts()!} eine zusätzliche
harte Schutzregel ein, die das Modell auf Vorschau- und Analyseinhalte
beschränkt. Die vier \passthrough{\lstinline!\_build\_*!}-Methoden
erzeugen die Bausteine unabhängig voneinander, sodass fehlende
Informationen (z. B. kein Spieler bei Phasenereignissen) den Prompt
nicht invalidieren, sondern den entsprechenden Block auslassen.

LLM-Parameter: \passthrough{\lstinline!temperature=0.3!} für konsistente
Ausgaben, \passthrough{\lstinline!max\_tokens!} aus
\passthrough{\lstinline!LLM\_MAX\_TOKENS!} (konfigurierbar). Bei
Rate-Limit-Fehlern werden bis zu
\passthrough{\lstinline!LLM\_RETRY\_ATTEMPTS!} (3) Versuche mit linearem
Backoff (\passthrough{\lstinline!30 s / 60 s / 90 s!}) unternommen; bei
sonstigen transienten Fehlern greift exponentielles Backoff
(\passthrough{\lstinline!1 s / 2 s / 4 s!}).

\textbf{Normalisierung des Event-Typs} erfolgt vorab in
\passthrough{\lstinline!\_normalize\_event\_type()!}. Die
\passthrough{\lstinline!EVENT\_TYPE\_MAP!} übersetzt externe Strings
(Partner-API: \passthrough{\lstinline!PartnerGoal!}, Football-API:
\passthrough{\lstinline!Goal!}) auf interne Bezeichner
(\passthrough{\lstinline!goal!}). Nicht erkannte Typen fallen auf
\passthrough{\lstinline!'comment'!} zurück.

\begin{center}\rule{0.5\linewidth}{0.5pt}\end{center}

\subsubsection{5.2.7 WebSocket-Endpunkt für
Media-Queue}\label{websocket-endpunkt-fuxfcr-media-queue}

Der WebSocket-Endpunkt \passthrough{\lstinline!/ws/media!} wird durch
einen dedizierten \passthrough{\lstinline!MediaConnectionManager!}
verwaltet:

\begin{lstlisting}[language=Python]
class MediaConnectionManager:
    def __init__(self) -> None:
        self.active_connections: list[WebSocket] = []

    async def connect(self, websocket: WebSocket) -> None:
        await websocket.accept()
        self.active_connections.append(websocket)

    async def broadcast(self, data: dict) -> None:
        dead: list[WebSocket] = []
        for ws in self.active_connections:
            try:
                await ws.send_json(data)
            except (WebSocketDisconnect, RuntimeError):
                dead.append(ws)
        for ws in dead:
            self.disconnect(ws)
\end{lstlisting}

Wenn das Backend neue Media-Queue-Einträge über
\passthrough{\lstinline!POST /api/v1/media/incoming!} erhält, ruft der
Router anschließend
\passthrough{\lstinline!manager.broadcast(\{'type': 'new\_media', 'items': [...]\})!}
auf. Alle verbundenen Clients erhalten das Update in Echtzeit. Der
Manager hält Verbindungen in einer einfachen In-Memory-Liste, was für
die aktuelle Nutzung (wenige gleichzeitige Redakteur-Clients)
ausreichend ist; bei Multi-Prozess-Deployments wäre ein externes
Pub/Sub-System (z. B. Redis) erforderlich.

Die persistente Seite der Media-Queue wird durch das
\passthrough{\lstinline!MediaQueueStatus!}-Enum gesteuert, das zwei
Zustände unterscheidet: \passthrough{\lstinline!pending!} (Medien sind
eingegangen und warten auf redaktionelle Zuordnung) und
\passthrough{\lstinline!published!} (Medien wurden einem Ticker-Eintrag
zugewiesen). Das Frontend zeigt im
\passthrough{\lstinline!MediaPickerPanel!} ausschließlich
\passthrough{\lstinline!pending!}-Einträge an; sobald ein Redakteur ein
Bild einem Ticker-Eintrag anhängt, wird der Status per
\passthrough{\lstinline!PATCH /api/v1/media/queue/\{id\}!} auf
\passthrough{\lstinline!published!} gesetzt. Dieses einfache
Zwei-Zustands-Modell verhindert, dass einmal zugeordnete Medien in der
Auswahl erneut auftauchen. Der korrespondierende Frontend-Hook
\passthrough{\lstinline!useMediaWebSocket!}, der die
Verbindungslebensdauer, das exponentielle Reconnect-Backoff und das
Nachrichten-Handling auf Client-Seite übernimmt, ist in Abschnitt 5.4.5
beschrieben.

\begin{center}\rule{0.5\linewidth}{0.5pt}\end{center}

\subsection{5.3 KI-Generierungspipeline}\label{ki-generierungspipeline}

\subsubsection{5.3.1 Ablaufdiagramm}\label{ablaufdiagramm}

Die folgende Darstellung zeigt den vollständigen Ablauf von einem neuen
Live-Ereignis bis zum publizierten Ticker-Eintrag:

\begin{lstlisting}
sequenceDiagram
    participant n8n as n8n Workflow
    participant API as FastAPI Backend
    participant DB as PostgreSQL
    participant LLM as LLM Provider
    participant FE as React Frontend

    n8n->>API: POST /api/v1/events (Ereignis-Import)
    API->>DB: INSERT INTO events
    n8n->>API: POST /api/v1/ticker/generate/{event_id}
    API->>DB: SELECT ticker_entry WHERE event_id (Dedup)
    alt Eintrag existiert bereits
        API-->>n8n: 200 OK (existierender Eintrag)
    else Neues Ereignis
        API->>DB: SELECT match + teams (joinedload)
        API->>DB: SELECT style_references (Few-Shot)
        API->>LLM: Prompt (Fakten + Kontext + Stilreferenzen)
        LLM-->>API: Generierter Text + Modellname
        API->>DB: INSERT ticker_entry (status=draft|published)
        API-->>n8n: 201 Created
    end
    FE->>API: GET /api/v1/ticker/{match_id} (Polling 5s)
    API-->>FE: TickerEntry[] (inkl. neuer Entwurf)
    alt Modus=coop
        FE-->>FE: Entwurf in Draft-Queue anzeigen
        note over FE: Redakteur: TAB (accept) / ESC (reject)
        FE->>API: PATCH /api/v1/ticker/{id} {status:'published'}
    else Modus=auto
        note over DB: status=published gesetzt bei Generierung
    end
\end{lstlisting}

\begin{center}\rule{0.5\linewidth}{0.5pt}\end{center}

\subsubsection{5.3.2 Deduplizierung}\label{deduplizierung}

Ein zentrales Qualitätsmerkmal der Pipeline ist die Deduplizierung: Der
Generierungsendpunkt prüft vor jedem LLM-Aufruf, ob bereits ein
Ticker-Eintrag für die \passthrough{\lstinline!event\_id!} existiert:

\begin{lstlisting}[language=Python]
existing = ticker_repo.get_by_event(event_id)
if existing:
    return existing  # Idempotenter Rückgabepfad ohne LLM-Aufruf
\end{lstlisting}

Dies macht den Endpunkt idempotent gegenüber Mehrfachaufrufen --- ein
wichtiges Merkmal, da n8n-Workflows bei Fehlern automatisch wiederholen.
Ohne diese Prüfung würde jeder Retry einen zusätzlichen Ticker-Eintrag
mit unterschiedlichem (nicht-deterministischem) Text erzeugen.

\begin{center}\rule{0.5\linewidth}{0.5pt}\end{center}

\subsubsection{5.3.3 Statusentscheidung und
Modus-Steuerung}\label{statusentscheidung-und-modus-steuerung}

Der initiale Status des erzeugten Ticker-Eintrags wird über den
Aufrufparameter \passthrough{\lstinline!auto\_publish!} gesteuert, den
n8n aus dem \passthrough{\lstinline!ticker\_mode!} des Spiels ableitet
(vgl. Kap. 4.3.3 für die konzeptionelle Beschreibung der drei
Betriebsmodi):

\begin{lstlisting}[language=Python]
status=TickerStatus.published if data.auto_publish else TickerStatus.draft
\end{lstlisting}

Damit wird die Modus-Logik zu einem reinen Aufrufparameter des Backends
--- die Entscheidungsinstanz ist n8n auf Basis des in der Datenbank
gespeicherten \passthrough{\lstinline!ticker\_mode!}. Das Backend selbst
ist modus-agnostisch: Es erhält einen Booleschen Wert und setzt den
Status entsprechend.

Die drei Modi übersetzen sich auf Implementierungsebene in drei
verschiedene Pfade:

{\def\LTcaptype{none} % do not increment counter
\begin{longtable}[]{@{}
  >{\raggedright\arraybackslash}p{(\linewidth - 6\tabcolsep) * \real{0.1270}}
  >{\raggedright\arraybackslash}p{(\linewidth - 6\tabcolsep) * \real{0.2222}}
  >{\raggedright\arraybackslash}p{(\linewidth - 6\tabcolsep) * \real{0.1746}}
  >{\raggedright\arraybackslash}p{(\linewidth - 6\tabcolsep) * \real{0.4762}}@{}}
\toprule\noalign{}
\begin{minipage}[b]{\linewidth}\raggedright
Modus
\end{minipage} & \begin{minipage}[b]{\linewidth}\raggedright
\passthrough{\lstinline!auto\_publish!}
\end{minipage} & \begin{minipage}[b]{\linewidth}\raggedright
Ergebnis
\end{minipage} & \begin{minipage}[b]{\linewidth}\raggedright
Redakteurs-Eingriff
\end{minipage} \\
\midrule\noalign{}
\endhead
\bottomrule\noalign{}
\endlastfoot
\passthrough{\lstinline!auto!} & \passthrough{\lstinline!True!} &
\passthrough{\lstinline!published!} & --- \\
\passthrough{\lstinline!coop!} & \passthrough{\lstinline!False!} &
\passthrough{\lstinline!draft!} & TAB (accept) / ESC (reject) \\
\passthrough{\lstinline!manual!} & --- & --- & Slash-Command →
\passthrough{\lstinline!POST /manual!} \\
\end{longtable}
}

Im \passthrough{\lstinline!manual!}-Modus wird der Generierungs-Endpunkt
gar nicht aufgerufen: n8n importiert zwar weiterhin Ereignisse, triggert
aber keinen LLM-Aufruf. Ticker-Einträge entstehen ausschließlich über
den \passthrough{\lstinline!POST /api/v1/ticker/manual!}-Endpunkt, den
das Frontend aus dem Slash-Command-Parser heraus bedient. Der Eintrag
wird dort direkt mit \passthrough{\lstinline!status=published!}
angelegt, da die Redakteurin den Text selbst verfasst hat und keine
weitere Freigabe erforderlich ist.

Im \passthrough{\lstinline!coop!}-Modus liegt die gesamte UX-Last auf
der Draft-Queue: Das Frontend zeigt neue
\passthrough{\lstinline!draft!}-Einträge mit
\passthrough{\lstinline!TAB!}/\passthrough{\lstinline!ESC!}-Shortcuts an
(vgl. Abschnitt 5.4.6), und die Entscheidungs-Latenz beträgt je nach
Komplexität des Entwurfs 15--30 Sekunden. Dieses Design stellt sicher,
dass keine KI-generierten Texte ohne menschliche Prüfung publiziert
werden --- eine Entscheidung, die direkt aus dem journalistischen
Qualitätsanspruch in Kapitel 2.1 folgt.

\begin{center}\rule{0.5\linewidth}{0.5pt}\end{center}

\subsubsection{5.3.4 Zustandsautomat der
Ticker-Einträge}\label{zustandsautomat-der-ticker-eintruxe4ge}

Das konzeptionelle Statusmodell (\passthrough{\lstinline!draft!} →
\passthrough{\lstinline!published!} /
\passthrough{\lstinline!rejected!}, Undo-Übergang
\passthrough{\lstinline!published → draft!}) ist vollständig in
Abschnitt 4.3.2 beschrieben. Implementierungsseitig sind zwei Details
hervorzuheben:

Der Übergang \passthrough{\lstinline!published → draft!} (Undo) ist im
Frontend als \textbf{Toast-Aktion} realisiert: Nach jeder Publikation
erscheint für 5 Sekunden ein Widerruf-Button. Über
\passthrough{\lstinline!PATCH /api/v1/ticker/\{id\}!} mit
\passthrough{\lstinline!\{ 'status': 'draft' \}!} wird der Eintrag
zurückgesetzt und steht erneut zur Bearbeitung bereit.

\passthrough{\lstinline!rejected!}-Einträge werden nicht aus der
Datenbank gelöscht, sondern behalten ihren Status. Diese Entscheidung
sichert die Reproduzierbarkeit der Evaluation: Alle je generierten
Einträge --- auch verworfene --- sind für spätere Qualitätsanalysen
auswertbar.

\begin{center}\rule{0.5\linewidth}{0.5pt}\end{center}

\subsection{5.4
Frontend-Implementierung}\label{frontend-implementierung}

\subsubsection{5.4.1 Komponentenhierarchie}\label{komponentenhierarchie}

Das Frontend ist als React-Anwendung mit TypeScript organisiert. Die
Hauptansicht (\passthrough{\lstinline!LiveTicker!}) gliedert sich in
drei Panel-Komponenten, die von gemeinsam genutzten Hooks und Contexts
gespeist werden:

\begin{lstlisting}
graph TD
    App --> LiveTicker
    LiveTicker --> |TickerModeContext| LeftPanel
    LiveTicker --> |TickerModeContext| CenterPanel
    LiveTicker --> |TickerModeContext| RightPanel
    LiveTicker --> |TickerDataContext| LeftPanel
    LiveTicker --> |TickerDataContext| CenterPanel
    LiveTicker --> |TickerDataContext| RightPanel
    LiveTicker --> |TickerActionsContext| CenterPanel

    CenterPanel --> EventCard
    CenterPanel --> EntryEditor
    CenterPanel --> AIDraft
    CenterPanel --> SummarySection
    CenterPanel --> MediaPickerPanel
    CenterPanel --> YouTubePanel
    CenterPanel --> TwitterPanel
    CenterPanel --> InstagramPanel

    LeftPanel --> PublishedEntry
    RightPanel --> Collapsible['Collapsible (Statistiken, Aufstellung, Spieler)']
    RightPanel --> FormationColumn
    RightPanel --> StatRow

    EventCard --> |onGenerate| TickerActionsContext
    AIDraft --> |onPublished/onDraftActive| TickerActionsContext
    EntryEditor --> |onManualPublish| TickerActionsContext
\end{lstlisting}

Die drei Panels sind physisch getrennte React-Komponenten, kommunizieren
aber ausschließlich über Context-Provider --- kein Prop-Drilling über
Komponentengrenzen hinweg.

\begin{center}\rule{0.5\linewidth}{0.5pt}\end{center}

\subsubsection{5.4.2 React
Context-Architektur}\label{react-context-architektur}

Das Frontend verwendet drei spezialisierte Contexts, die in
\passthrough{\lstinline!LiveTicker!} zusammengeführt werden. Jeder
Context hat eine klar definierte Verantwortung:

\textbf{\passthrough{\lstinline!TickerModeContext!}} --- Modus-State und
Tastatur-Aktionen:

\begin{lstlisting}
interface TickerModeContextValue {
  mode: TickerMode; // 'auto' | 'coop' | 'manual'
  setMode: (mode: TickerMode) => void;
  acceptDraft: () => Promise<void>; // TAB-Shortcut
  rejectDraft: () => Promise<void>; // ESC-Shortcut
}
\end{lstlisting}

\textbf{\passthrough{\lstinline!TickerDataContext!}} --- Match-Daten und
Reload-Funktionen:

\begin{lstlisting}
interface TickerDataContextValue {
  match: Match | null;
  events: MatchEvent[];
  tickerTexts: TickerEntry[];
  prematch: TickerEntry[];
  lineups: LineupEntry[];
  matchStats: MatchStat[];
  players: Player[];
  playerStats: PlayerStat[];
  injuries: InjuryGroup[];
  reload: ReloadFunctions;
  generatingId: number | null;
}
\end{lstlisting}

\textbf{\passthrough{\lstinline!TickerActionsContext!}} --- Callbacks
für redaktionelle Aktionen:

\begin{lstlisting}
interface TickerActionsContextValue {
  onGenerate: (eventId: number, style: TickerStyle) => Promise<void>;
  onManualPublish: (text, icon?, minute?, phase?, rawInput?) => Promise<void>;
  onDraftActive: (id: number, text: string) => void;
  onPublished: (id: number, text: string, isManual?: boolean) => void;
  onEditEntry: (id: number, text: string) => Promise<void>;
  onDeleteEntry: (id: number) => Promise<void>;
  retractedText: string | null;
  clearRetractedText: () => void;
}
\end{lstlisting}

Die Trennung von Daten- und Aktions-Context hat einen konkreten
Performance-Vorteil: Komponenten, die nur Actions konsumieren (z. B.
\passthrough{\lstinline!EntryEditor!}), re-rendern nicht bei Änderungen
an \passthrough{\lstinline!tickerTexts!}. Dies ist bei einem Live-Ticker
mit 5-Sekunden-Polling relevant.

Alle Contexts exportieren zusätzlich einen Hook
(\passthrough{\lstinline!useTickerModeContext!},
\passthrough{\lstinline!useTickerDataContext!},
\passthrough{\lstinline!useTickerActionsContext!}), der bei Verwendung
außerhalb des Providers einen klaren Fehler wirft --- ein
Entwicklungssicherheitsnetz.

\begin{center}\rule{0.5\linewidth}{0.5pt}\end{center}

\subsubsection{5.4.3 Hook-Architektur}\label{hook-architektur}

Die Zustandslogik ist in zwei Ebenen von Hooks organisiert: gemeinsame
Hooks unter \passthrough{\lstinline!src/hooks/!} und feature-spezifische
Hooks unter \passthrough{\lstinline!components/LiveTicker/hooks/!}.

Die wichtigsten Hooks und ihre Verantwortlichkeiten:

{\def\LTcaptype{none} % do not increment counter
\begin{longtable}[]{@{}
  >{\raggedright\arraybackslash}p{(\linewidth - 4\tabcolsep) * \real{0.2043}}
  >{\raggedright\arraybackslash}p{(\linewidth - 4\tabcolsep) * \real{0.0968}}
  >{\raggedright\arraybackslash}p{(\linewidth - 4\tabcolsep) * \real{0.6989}}@{}}
\toprule\noalign{}
\begin{minipage}[b]{\linewidth}\raggedright
Hook
\end{minipage} & \begin{minipage}[b]{\linewidth}\raggedright
Paket
\end{minipage} & \begin{minipage}[b]{\linewidth}\raggedright
Verantwortlichkeit
\end{minipage} \\
\midrule\noalign{}
\endhead
\bottomrule\noalign{}
\endlastfoot
\passthrough{\lstinline!useNavigation!} & hooks/ &
Land→Team→Wettbewerb→Spieltag→Spiel Navigationskette \\
\passthrough{\lstinline!useMatchData!} & hooks/ & Aggregiert die drei
Polling-Hooks und koordiniert Reload-Zyklen \\
\passthrough{\lstinline!useMatchCore!} & hooks/ & Polling der
Match-Metadaten (Teams, Status, Spielstand, Minute) \\
\passthrough{\lstinline!useMatchEvents!} & hooks/ & Polling der
Spielereignisse (Tore, Karten, Auswechslungen) \\
\passthrough{\lstinline!useMatchTicker!} & hooks/ & Polling der
Ticker-Einträge (Entwürfe, publizierte Texte) \\
\passthrough{\lstinline!useMatchTriggers!} & hooks/ & n8n-Webhooks nach
Matchauswahl und Statuswechseln \\
\passthrough{\lstinline!useTickerMode!} & hooks/ &
\passthrough{\lstinline!mode!}-State, PATCH zum Backend,
Keyboard-Shortcuts \\
\passthrough{\lstinline!useMediaWebSocket!} & hooks/ &
WebSocket-Verbindung mit exponentiellem Reconnect-Backoff \\
\passthrough{\lstinline!usePanelResize!} & hooks/ & Drag-Resize zwischen
Center- und Right-Panel \\
\passthrough{\lstinline!useApiStatus!} & hooks/ & Health-Check gegen
Backend (alle 30 s) \\
\passthrough{\lstinline!useTicker!} & LT/hooks/ & Ticker-Actions,
Draft-Queue, Publish-Toast (Undo) \\
\passthrough{\lstinline!useEventDraft!} & LT/hooks/ & Entwurfsansicht,
Accept/Reject-Logik \\
\passthrough{\lstinline!useAutoPublisher!} & LT/hooks/ & Automatisches
Publizieren im Auto-Modus \\
\passthrough{\lstinline!useRightPanelData!} & LT/hooks/ & Aufbereitung
von Lineups, Stats, Spielerinfo für das rechte Panel \\
\end{longtable}
}

\passthrough{\lstinline!useMatchCore!},
\passthrough{\lstinline!useMatchEvents!} und
\passthrough{\lstinline!useMatchTicker!} betreiben jeweils einen eigenen
\passthrough{\lstinline!setInterval!}-Zyklus, der bei Mount startet und
bei Unmount gestoppt wird. Die HTTP-Kommunikation erfolgt über
\textbf{Axios} als HTTP-Client. \passthrough{\lstinline!useMatchData!}
aggregiert die drei Hooks und stellt die Daten über
\passthrough{\lstinline!TickerDataContext!} bereit.

\passthrough{\lstinline!useMatchData!} steuert das Polling-Intervall
über die Hilfsfunktion
\passthrough{\lstinline!resolvePollingInterval()!}:

\begin{lstlisting}
// src/utils/resolvePollingInterval.ts
export function resolvePollingInterval(matchState: string | null): number {
  if (matchState === 'Live' || matchState === 'FullTime') return POLL_EVENTS_MS; // 5000
  if (matchState == null) return POLL_EVENTS_MS; // 5000
  return POLL_PREMATCH_MS; // 5000
}
\end{lstlisting}

Aktuell sind beide Konstanten
(\passthrough{\lstinline!POLL\_EVENTS\_MS!},
\passthrough{\lstinline!POLL\_PREMATCH\_MS!}) auf 5000 ms gesetzt,
sodass alle Zustände mit einem einheitlichen 5-Sekunden-Intervall
gepollt werden. Die Unterscheidung ist als Erweiterungspunkt vorgesehen,
um bei Bedarf ruhigere Polling-Intervalle für inaktive Spiele
einzuführen.

\begin{center}\rule{0.5\linewidth}{0.5pt}\end{center}

\subsubsection{\texorpdfstring{5.4.4 Slash-Command-Parser
(\texttt{parseCommand.ts})}{5.4.4 Slash-Command-Parser (parseCommand.ts)}}\label{slash-command-parser-parsecommand.ts}

Der Slash-Command-Parser ermöglicht schnelle manuelle Texteingabe ohne
Maus-Interaktion. Er übersetzt kurze Kommandos in formatierte
Ticker-Texte und liefert Metadaten (Icon, Phase, Minute) für die
Publikation.

Die Funktion
\passthrough{\lstinline!parseCommand(input, currentMinute)!} gibt ein
typisiertes \passthrough{\lstinline!ParseResult!} zurück:

\begin{lstlisting}
export interface ParseResult {
  type: string;
  formatted: string;
  warnings: string[];
  isValid: boolean;
  meta: {
    icon: string;
    phase: string | null;
    minute: number | null;
  };
}
\end{lstlisting}

\textbf{Phasen-Commands} (11 Einträge in
\passthrough{\lstinline!PHASE\_CMDS!}) erzeugen vordefinierte Texte mit
optionaler Minutenangabe. Die folgende Tabelle zeigt eine Auswahl der
häufigsten Einträge (7 von 11):

{\def\LTcaptype{none} % do not increment counter
\begin{longtable}[]{@{}llll@{}}
\toprule\noalign{}
Command & Ausgabe & Phase & Icon \\
\midrule\noalign{}
\endhead
\bottomrule\noalign{}
\endlastfoot
\passthrough{\lstinline!/prematch!} & ``Vorbericht'' &
\passthrough{\lstinline!Before!} & 📣 \\
\passthrough{\lstinline!/anpfiff!} & ``Anpfiff!'' &
\passthrough{\lstinline!FirstHalf!} & 📣 \\
\passthrough{\lstinline!/hz!} & ``Halbzeit!'' &
\passthrough{\lstinline!FirstHalfBreak!} & 🔔 \\
\passthrough{\lstinline!/2hz!} & ``Anstoß zur 2. Halbzeit'' &
\passthrough{\lstinline!SecondHalf!} & 📣 \\
\passthrough{\lstinline!/vz1!} & ``Beginn der Verlängerung'' &
\passthrough{\lstinline!ExtraFirstHalf!} & 📣 \\
\passthrough{\lstinline!/elfmeter!} & ``Elfmeterschießen beginnt'' &
\passthrough{\lstinline!PenaltyShootout!} & 🥅 \\
\passthrough{\lstinline!/abpfiff!} & ``Abpfiff!'' &
\passthrough{\lstinline!After!} & 📣 \\
\end{longtable}
}

\textbf{Event-Commands} (12 Einträge in
\passthrough{\lstinline!CMD\_MAP!}, 8 Typen nach Expansion) erzeugen
strukturierte Texte mit Validierungshinweisen:

\begin{lstlisting}
/g Mustermann EF  → 'TOR — Mustermann (EF)'  + icon ⚽
/gelb Müller EF   → 'Gelb — Müller (EF)'     + icon 🟨
/rot Huber Bayern → 'Rote Karte — Huber (Bayern)' + icon 🟥
/s EinRaus AusRein EF → 'Wechsel — EinRaus ↔ AusRein (EF)' + icon 🔄
/n Freistoß knapp über das Tor → 'Freistoß knapp über das Tor'
\end{lstlisting}

Für unvollständige Eingaben gibt \passthrough{\lstinline!warnings[]!}
gezielte Hinweise zurück (\passthrough{\lstinline!'Fehlend: Spieler'!}),
die im \passthrough{\lstinline!EntryEditor!} live angezeigt werden. Das
Frontend kann damit ohne Backend-Aufruf Validierungsfeedback geben.

Nicht erkannte Commands werden mit
\passthrough{\lstinline!isValid: false!} und einer Fehlermeldung
zurückgegeben (\passthrough{\lstinline!'Unbekannter Command: /xyz'!}),
statt einen Laufzeitfehler zu werfen.

\begin{center}\rule{0.5\linewidth}{0.5pt}\end{center}

\subsubsection{\texorpdfstring{5.4.5 WebSocket-Hook
(\texttt{useMediaWebSocket})}{5.4.5 WebSocket-Hook (useMediaWebSocket)}}\label{websocket-hook-usemediawebsocket}

Der \passthrough{\lstinline!useMediaWebSocket!}-Hook verwaltet die
Lebensdauer der WebSocket-Verbindung zum Backend. Die Kernentscheidung
ist das \textbf{exponentielle Reconnect-Backoff}:

\begin{lstlisting}
const delay = Math.min(
  BASE_DELAY_MS * 2 ** retryCountRef.current, // 1s → 2s → 4s → 8s → …
  MAX_DELAY_MS, // max 30s
);
\end{lstlisting}

Der \passthrough{\lstinline!retryCountRef!} zählt aufeinanderfolgende
Verbindungsabbrüche. Bei erfolgreicher Reconnection wird er auf 0
zurückgesetzt. \passthrough{\lstinline!BASE\_DELAY\_MS = 1000!},
\passthrough{\lstinline!MAX\_DELAY\_MS = 30\_000!}.

Ein \passthrough{\lstinline!mountedRef!} verhindert Zustandsänderungen
nach dem Unmounten der Komponente --- ein häufiges React-Problem bei
asynchronen Operationen:

\begin{lstlisting}
ws.onclose = () => {
  if (!mountedRef.current) return; // Kein setState nach Unmount
  // … retry-Logik
};
\end{lstlisting}

Der Cleanup im \passthrough{\lstinline!useEffect!} setzt
\passthrough{\lstinline!ws.onclose = null!}, bevor die WebSocket
geschlossen wird. Ohne diesen Schritt würde das
\passthrough{\lstinline!close!}-Event des kontrollierten Schließens
einen unbeabsichtigten Reconnect-Versuch auslösen.

\begin{center}\rule{0.5\linewidth}{0.5pt}\end{center}

\subsubsection{5.4.6 Modusimplementierung im
Frontend}\label{modusimplementierung-im-frontend}

Der aktive Modus (\passthrough{\lstinline!auto!} /
\passthrough{\lstinline!coop!} / \passthrough{\lstinline!manual!})
steuert das Verhalten mehrerer Frontend-Komponenten (vgl. Kap. 4.6.4 für
die konzeptionelle Einordnung der Moduslogik und Interaktionsdesign).
Die Implementierung verteilt sich auf drei Ebenen:

\textbf{Persistenz und Synchronisation} ---
\passthrough{\lstinline!useTickerMode!} initialisiert den Modus lokal
mit \passthrough{\lstinline!coop!} als Standardwert und überschreibt
diesen beim Laden eines Spiels mit dem in der Datenbank gespeicherten
\passthrough{\lstinline!ticker\_mode!}. Änderungen werden über
\passthrough{\lstinline!PATCH /matches/\{id\}/ticker-mode!} ans Backend
zurückgeschrieben. Lokale Optimistic-Updates vermeiden wahrnehmbare
Latenz beim Umschalten.

\textbf{Keyboard-Shortcuts} --- Im Co-op-Modus registriert
\passthrough{\lstinline!useTickerMode!} die Shortcuts
\passthrough{\lstinline!TAB!} (acceptDraft) und
\passthrough{\lstinline!ESC!} (rejectDraft) direkt als
Keyboard-Event-Listener. Die Shortcut-Registrierung ist an den
Modus-Zustand gebunden und wird bei Moduswechsel automatisch
aktualisiert.

\textbf{\passthrough{\lstinline!useAutoPublisher!}} --- Im Auto-Modus
überwacht dieser Hook die \passthrough{\lstinline!tickerTexts!} auf neue
Draft-Einträge (\passthrough{\lstinline!status='draft'!}) und publiziert
sie automatisch. Dies ist als Frontend-Fallback gedacht: Primär setzt
n8n mit \passthrough{\lstinline!auto\_publish=True!} den Status direkt
beim Generieren; \passthrough{\lstinline!useAutoPublisher!} fängt Fälle
auf, in denen das nicht geschehen ist.

\begin{center}\rule{0.5\linewidth}{0.5pt}\end{center}

\subsubsection{5.4.7 TypeScript-Migration}\label{typescript-migration}

Das Frontend wurde im Verlauf des Projekts von reinem JavaScript auf
TypeScript migriert. Ziel war nicht die vollständige Auflösung aller
\passthrough{\lstinline!any!}-Typen, sondern eine pragmatische
Typisierung der systemkritischen Pfade mit messbaren
Qualitätskennzahlen.

\textbf{Migrationsstrategie} --- Die Migration folgte einem
schrittweisen Ansatz: Zuerst wurden alle \passthrough{\lstinline!.js!}-
und \passthrough{\lstinline!.jsx!}-Dateien in
\passthrough{\lstinline!.ts!} und \passthrough{\lstinline!.tsx!}
umbenannt, dann \passthrough{\lstinline!tsconfig.json!} konfiguriert
(\passthrough{\lstinline!target: es2015!},
\passthrough{\lstinline!module: esnext!},
\passthrough{\lstinline!jsx: react-jsx!},
\passthrough{\lstinline!moduleResolution: node!}). Anschließend wurden
TypeScript-Fehler von innen nach außen behoben --- beginnend mit dem
Typ-System (\passthrough{\lstinline!src/types/index.ts!}) über Contexts
bis zu den Komponenten. Der \passthrough{\lstinline!strict!}-Modus ist
derzeit deaktiviert (\passthrough{\lstinline!'strict': false!}), um die
inkrementelle Migration ohne blockierende Fehlereskalation zu
ermöglichen.

\textbf{Zentrale Typ-Definitionen} ---
\passthrough{\lstinline!src/types/index.ts!} definiert die gemeinsamen
Domain-Interfaces:

\begin{lstlisting}
export type TickerMode = 'auto' | 'coop' | 'manual';
export type TickerStyle = 'neutral' | 'euphorisch' | 'kritisch';
export type MatchPhase = 'Before' | 'FirstHalf' | 'FirstHalfBreak' | 'SecondHalf' | ...;

export interface TickerEntry {
  id: number;
  match_id: number;
  event_id?: number | null;
  text: string;
  style?: string | null;
  status: 'draft' | 'published' | 'rejected';
  source?: 'ai' | 'manual';
  minute?: number | null;
  phase?: string | null;
  icon?: string | null;
  image_url?: string | null;
  video_url?: string | null;
  llm_model?: string | null;
}
\end{lstlisting}

\textbf{Ergebnis} --- Nach der Migration erzielt das Projekt \textbf{0
TypeScript-Compilerfehler} und eine \textbf{type-coverage von 95,84 \%}
(gemessen mit \passthrough{\lstinline!type-coverage --strict!}). Die
detaillierte Darstellung der Migrationsergebnisse und die Analyse der
verbleibenden untypisierten Stellen folgt in Abschnitt 6.2.5.

\begin{center}\rule{0.5\linewidth}{0.5pt}\end{center}

\subsection{5.5 n8n
Workflow-Implementierung}\label{n8n-workflow-implementierung}

Die n8n-Workflows bilden die Orchestrierungsschicht zwischen externen
Datenquellen, dem FastAPI-Backend und dem LLM-Dienst. Alle 15
Workflow-Dateien liegen als versionierte JSON-Exporte im
Projektverzeichnis \passthrough{\lstinline!n8n/!} vor und können direkt
in eine n8n-Instanz importiert werden. Die Workflows gliedern sich in
vier funktionale Gruppen:

\textbf{A) Stammdaten und Matchstruktur} --- Die Workflows
\passthrough{\lstinline!01\_import\_countries!},
\passthrough{\lstinline!02\_import\_teams\_by\_country!},
\passthrough{\lstinline!03\_import\_competitions\_and\_matches!} und
\passthrough{\lstinline!04\_import\_matches!} laden Länder, Teams,
Wettbewerbe und Spielpläne aus API-Football und persistieren diese via
Upsert-Strategien. Externe IDs dienen als stabile Zuordnungsschlüssel
für konfliktfreie Aktualisierungen.

\textbf{B) Spielbezogene Detaildaten} ---
\passthrough{\lstinline!04\_import\_lineups!},
\passthrough{\lstinline!05\_import\_match\_statistics!},
\passthrough{\lstinline!06\_import\_player\_statistics!} und
\passthrough{\lstinline!07\_import\_prematch!} importieren
Aufstellungen, Match- und Spielerstatistiken matchbezogen. Der
Prematch-Workflow erzeugt zusätzlich strukturierte Kontextdaten
(Verletzungen, Head-to-Head, Tabellenstand) als synthetische Events, die
dem LLM als Promptkontext dienen.

\textbf{C) KI-Generierung} ---
\passthrough{\lstinline!09\_events\_llm\_workflow!},
\passthrough{\lstinline!13\_Halftime\_aftertime!} und
\passthrough{\lstinline!14\_Game\_ANpfiff\_ABpfiff!} importieren
Live-Ereignisse, erzeugen synthetische Phasenereignisse und delegieren
die Textgenerierung an das Backend. Der Zusammenfassungs-Workflow
(\passthrough{\lstinline!13!}) ruft den LLM-Provider direkt auf, da
Halbzeit-/Abpfiff-Zusammenfassungen einen umfangreicheren Prompt mit
Statistiken erfordern.

\textbf{D) Medien und Social Media} ---
\passthrough{\lstinline!08\_scoreplay\_media\_workflow!},
\passthrough{\lstinline!10\_Twitter!},
\passthrough{\lstinline!11\_youtube!} und
\passthrough{\lstinline!12\_insta!} importieren Medien- und
Social-Media-Inhalte. Der Media-Workflow übergibt Bilder an das Backend,
das diese über WebSocket verteilt; die Social-Media-Workflows speichern
Clips in der Datenbank. Diese Workflows dienen der Ingestion, nicht dem
Publishing.

Dieser Abschnitt dokumentiert die Implementierungsdetails der fünf
zentralen Workflows (Gruppen B--D).

\begin{center}\rule{0.5\linewidth}{0.5pt}\end{center}

\subsubsection{\texorpdfstring{5.5.1 Events-LLM-Workflow
(\texttt{09\_events\_llm\_workflow.json})}{5.5.1 Events-LLM-Workflow (09\_events\_llm\_workflow.json)}}\label{events-llm-workflow-09_events_llm_workflow.json}

Workflow \passthrough{\lstinline!09!} ist der kritischste im System: Er
importiert Live-Ereignisse via Webhook
(\passthrough{\lstinline!POST /Events!}), persistiert sie per UPSERT
(\passthrough{\lstinline!ON CONFLICT (source\_id) DO NOTHING!}) und
triggert die KI-Generierung für jeden neuen Event. Eine initiale
SQL-Abfrage bestimmt anhand der Teamzugehörigkeit, ob die Instanz
\passthrough{\lstinline!ef\_whitelabel!} oder
\passthrough{\lstinline!generic!} und der Stil
\passthrough{\lstinline!euphorisch!} oder
\passthrough{\lstinline!neutral!} ist;
\passthrough{\lstinline!auto\_publish!} wird direkt aus
\passthrough{\lstinline!ticker\_mode!} des Spiels gesetzt. Nur Events
mit erfolgreicher Datenbankzeile (zurückgegebene
\passthrough{\lstinline!id!}, kein
\passthrough{\lstinline!DO NOTHING!}-Konflikt) passieren den
Filter-Knoten und triggern den Backend-Endpunkt. Ein zusätzlicher
EF-spezifischer Zweig liest Spielerdaten von
\passthrough{\lstinline!profis.eintracht.de!}, baut S3-Video-URLs aus
neunstellig aufgefüllten Spieler-IDs auf und persistiert Torjubel-Videos
als Ticker-Einträge.

\begin{center}\rule{0.5\linewidth}{0.5pt}\end{center}

\subsubsection{\texorpdfstring{5.5.2 Prematch-Import-Workflow
(\texttt{07\_import\_prematch.json})}{5.5.2 Prematch-Import-Workflow (07\_import\_prematch.json)}}\label{prematch-import-workflow-07_import_prematch.json}

Workflow \passthrough{\lstinline!07!} baut den Vorberichtskontext aus
fünf parallelen Football-API-Abfragen auf (Verletzungen, Head-to-Head,
Teamstatistiken Heim/Gast, Tabellenstand) und persistiert die Ergebnisse
als \passthrough{\lstinline!synthetic\_events!} mit idempotenter
\passthrough{\lstinline!ON CONFLICT (match\_id, type) DO UPDATE!}-Strategie.
Ein LLM-Trigger wird nur ausgelöst, wenn noch kein nicht-verworfener
Ticker-Eintrag für das jeweilige synthetische Event existiert
(\passthrough{\lstinline!WHERE te.id IS NULL!}) --- bereits generierte
und eventuell publizierte Texte werden damit nicht überschrieben.

\begin{center}\rule{0.5\linewidth}{0.5pt}\end{center}

\subsubsection{\texorpdfstring{5.5.3 Matchphasen-Workflow
(\texttt{14\_Game\_ANpfiff\_ABpfiff.json})}{5.5.3 Matchphasen-Workflow (14\_Game\_ANpfiff\_ABpfiff.json)}}\label{matchphasen-workflow-14_game_anpfiff_abpfiff.json}

Workflow \passthrough{\lstinline!14!} verarbeitet
Spielzustands-Übergänge (Anpfiff, Halbzeit, Abpfiff etc.) via Webhook
(\passthrough{\lstinline!POST /match-status!}). Ein JavaScript-Knoten
validiert Zustandsübergänge gegen eine Matrix erlaubter Vorzustände ---
ungültige Übergänge (z. B. direkt von \passthrough{\lstinline!PreMatch!}
zu \passthrough{\lstinline!2H!}) werden zurückgewiesen. Für
Vollzeit-Ereignisse (\passthrough{\lstinline!FT!},
\passthrough{\lstinline!AET!}, \passthrough{\lstinline!PEN!}) generiert
der Workflow die gesamte Phasensequenz rückwirkend: Ein
\passthrough{\lstinline!FT!}-Signal erzeugt vier synthetische Events
(Anstoß, Halbzeit, 2. Halbzeit, Abpfiff). Das zentrale SQL-Statement
kombiniert Match-Update, Event-Insert und eine
\passthrough{\lstinline!demote!}-CTE in einer einzigen Transaktion ---
Letztere stuft bereits publizierte Phasen-Texte bei Re-Triggers auf
\passthrough{\lstinline!draft!} zurück, sodass die Redaktion sie erneut
prüfen kann.

\begin{center}\rule{0.5\linewidth}{0.5pt}\end{center}

\subsubsection{\texorpdfstring{5.5.4 Halbzeit/Abpfiff-Zusammenfassung
(\texttt{13\_Halftime\_aftertime.json})}{5.5.4 Halbzeit/Abpfiff-Zusammenfassung (13\_Halftime\_aftertime.json)}}\label{halbzeitabpfiff-zusammenfassung-13_halftime_aftertime.json}

Workflow \passthrough{\lstinline!13!} erzeugt narrative
Zusammenfassungen für Halbzeit und Abpfiff. Im Gegensatz zu
\passthrough{\lstinline!09!} ruft er OpenRouter
(\passthrough{\lstinline!google/gemini-2.0-flash-lite-001!})
\textbf{direkt} auf, da Zusammenfassungen keinen Few-Shot-Stil aus der
Datenbank benötigen, sondern einen umfangreicheren Prompt mit parallel
geladenen Statistiken (Ballbesitz, Schüsse, Pässe, Spieler-Ratings)
erfordern. Das Ergebnis wird über
\passthrough{\lstinline!POST /api/v1/ticker/manual!} als Ticker-Eintrag
gespeichert --- je nach \passthrough{\lstinline!ticker\_mode!} direkt
als \passthrough{\lstinline!published!} oder als
\passthrough{\lstinline!draft!}.

\begin{center}\rule{0.5\linewidth}{0.5pt}\end{center}

\subsubsection{\texorpdfstring{5.5.5 ScorePlay-Medien-Workflow
(\texttt{08\_scoreplay\_media\_workflow.json})}{5.5.5 ScorePlay-Medien-Workflow (08\_scoreplay\_media\_workflow.json)}}\label{scoreplay-medien-workflow-08_scoreplay_media_workflow.json}

Workflow \passthrough{\lstinline!08!} sucht Medien-Assets bei ScorePlay
für Spieler eines Torereignisses. Die Spieler-Suche
(\passthrough{\lstinline!GET /v1/tag/search!}) normalisiert Umlaute
(ä→a, ö→o, ü→u, ß→ss) und matcht gegen vier Namensfelder
(\passthrough{\lstinline!full\_name!},
\passthrough{\lstinline!first\_name!},
\passthrough{\lstinline!last\_name!},
\passthrough{\lstinline!ai\_name!}). Gefundene Thumbnail-, Compressed-
und Original-URLs werden an
\passthrough{\lstinline!POST /api/v1/media/incoming!} übertragen und vom
Backend per WebSocket an verbundene Frontend-Clients verteilt.

\begin{center}\rule{0.5\linewidth}{0.5pt}\end{center}

\subsubsection{5.5.6 Implementierungsübergreifende
Muster}\label{implementierungsuxfcbergreifende-muster}

Über alle Workflows hinweg lassen sich fünf Implementierungsmuster
identifizieren:

\textbf{1. Idempotente UPSERT-Strategie} --- Alle Insert-Operationen
verwenden \passthrough{\lstinline!ON CONFLICT DO UPDATE!} oder
\passthrough{\lstinline!DO NOTHING!}. Als Konfliktschlüssel dienen
externe IDs (\passthrough{\lstinline!source\_id!},
\passthrough{\lstinline!external\_id!}, \passthrough{\lstinline!vid!}),
zusammengesetzte Schlüssel (\passthrough{\lstinline!match\_id, type!})
oder Unique-Constraints auf Namensfeldern. Das macht alle Workflows bei
Mehrfachausführung sicher.

\textbf{2. Filter-vor-Trigger-Muster} --- Vor jedem LLM-Trigger steht
ein Filter-Knoten, der nur Rows mit einer \passthrough{\lstinline!id!}
(erfolgreicher DB-Insert) durchlässt. Events mit
\passthrough{\lstinline!DO NOTHING!}-Konflikten (bereits importiert)
erzeugen keinen neuen LLM-Aufruf.

\textbf{3. Statusbasierte Filterung} --- Bei synthetischen Events und
Ticker-Einträgen werden vorhandene Zustände berücksichtigt, um
redundante Generierungen zu vermeiden: Ein LLM-Aufruf wird nur
getriggert, wenn noch kein nicht-verworfener Ticker-Eintrag für das
Event existiert (vgl. die WHERE-Bedingungen in 5.5.2 und 5.5.3). Bei
Re-Triggern werden bereits publizierte Phasen-Texte auf
\passthrough{\lstinline!draft!} zurückgestuft
(\passthrough{\lstinline!demote!}-CTE in 5.5.3).

\textbf{4. Ticker-Mode-Propagation} ---
\passthrough{\lstinline!ticker\_mode!} wird in allen
Generierungsworkflows als \passthrough{\lstinline!auto\_publish!}-Flag
an das Backend weitergegeben (vgl. Kap. 5.3.3). Die Moduslogik sitzt
damit vollständig in der Datenbank; n8n liest sie bei jedem Aufruf neu
aus.

\textbf{5. Instanz-Routing} --- Die Unterscheidung zwischen
\passthrough{\lstinline!ef\_whitelabel!} und
\passthrough{\lstinline!generic!} erfolgt per
\passthrough{\lstinline!ILIKE '\%Frankfurt\%'!}-Abfrage auf Teamnamen.
Das Ergebnis steuert Stilwahl, Few-Shot-Auswahl im Backend und ob der
Torjubel-Video-Pfad aktiviert wird.

\begin{center}\rule{0.5\linewidth}{0.5pt}\end{center}

\subsection{5.6 Qualitätssicherung und
Tests}\label{qualituxe4tssicherung-und-tests}

Die Qualitätssicherung folgt dem Testpyramiden-Modell nach Cohn (2009)
mit drei Ebenen: Unit-Tests, Integrations-Tests und End-to-End-Tests.
Insgesamt umfasst die Testsuite \textbf{391 Tests} (187 Frontend, 198
Backend, 6 E2E), die alle grün durchlaufen.

\textbf{Frontend-Tests} --- Das Frontend verwendet Jest mit
\passthrough{\lstinline!@testing-library/react!} für 187 Tests in 15
Dateien. Die Testdateien sind kolokiert mit ihrem Quellcode:
Komponentendateien unter
\passthrough{\lstinline!components/LiveTicker/components/!}, Hook-Tests
unter \passthrough{\lstinline!hooks/!} und Utility-Tests unter
\passthrough{\lstinline!utils/!}. Besonders umfangreich getestet ist der
\passthrough{\lstinline!parseCommand!}-Parser (45 Testfälle), da er die
kritische manuelle Eingabeschicht bildet. Ein exemplarischer Testfall
illustriert das Prüfmuster:

\begin{lstlisting}
test('vollstaendiger Command ist valid', () => {
  const result = parseCommand('/g Muller EIN', 32);
  expect(result.isValid).toBe(true);
  expect(result.type).toBe('goal');
  expect(result.formatted).toBe('TOR -- Muller (EIN)');
  expect(result.meta.icon).toBe('⚽');
  expect(result.warnings).toHaveLength(0);
});
\end{lstlisting}

\textbf{Backend-Tests} --- Das Backend nutzt pytest mit FastAPI
\passthrough{\lstinline!TestClient!} und transaktionalem Rollback für
198 Tests bei 75 \% Statement-Coverage. Die Tests gliedern sich in
API-Integrations-Tests (\passthrough{\lstinline!test\_ticker\_api.py!},
\passthrough{\lstinline!test\_matches\_api.py!},
\passthrough{\lstinline!test\_events\_api.py!} u. a.), die Endpunkte
gegen eine PostgreSQL-Testdatenbank prüfen, sowie Unit-Tests für
Services (\passthrough{\lstinline!test\_llm\_service.py!},
\passthrough{\lstinline!test\_ticker\_service.py!}) und Repositories
(\passthrough{\lstinline!test\_ticker\_entry\_repository.py!}).
Gemeinsame Fixtures (Datenbankverbindung, Test-Match, Test-Events) sind
in \passthrough{\lstinline!conftest.py!} zentralisiert.

\textbf{End-to-End-Tests} --- Sechs Playwright-Tests validieren den
stabilen Kern der Browser-Anwendung: korrektes Rendern, Fehlerfreiheit
beim Laden und grundlegende Interaktionspunkte.

\textbf{Evaluations-Infrastruktur} --- Ergänzend zur Testsuite enthält
\passthrough{\lstinline!backend/app/utils/evaluation\_metrics.py!} sechs
statistische Hilfsfunktionen, die speziell für die Qualitätsevaluation
der KI-Textgenerierung implementiert wurden:

{\def\LTcaptype{none} % do not increment counter
\begin{longtable}[]{@{}
  >{\raggedright\arraybackslash}p{(\linewidth - 2\tabcolsep) * \real{0.2824}}
  >{\raggedright\arraybackslash}p{(\linewidth - 2\tabcolsep) * \real{0.7176}}@{}}
\toprule\noalign{}
\begin{minipage}[b]{\linewidth}\raggedright
Funktion
\end{minipage} & \begin{minipage}[b]{\linewidth}\raggedright
Aufgabe
\end{minipage} \\
\midrule\noalign{}
\endhead
\bottomrule\noalign{}
\endlastfoot
\passthrough{\lstinline!compute\_ttp\_stats(entries)!} & Berechnet
TTP-Statistiken (Mean, Median, Perzentile) aus gemessenen Latenzen \\
\passthrough{\lstinline!cliffs\_delta(group\_a, group\_b)!} &
Effektstärkemaß nach Cliff (1993) für ordinal-skalierte Vergleiche (z.
B. auto vs.~manual TTP) \\
\passthrough{\lstinline!bootstrap\_ci(data, n\_resamples)!} &
Bootstrap-Konfidenzintervall (95 \%) über nichtparametrisches
Resampling \\
\passthrough{\lstinline!cohens\_kappa(rater\_a, rater\_b)!} &
Inter-Rater-Übereinstimmungsmaß nach Cohen (1960) für die
Qualitätsbewertung \\
\passthrough{\lstinline!distribution\_summary(scores)!} & Deskriptive
Zusammenfassung (Mean, SD, Min, Max) für Bewertungslisten \\
\passthrough{\lstinline!aggregate\_quality\_by\_group(entries)!} &
Aggregiert Qualitätsscores nach Gruppe (Stil, Ereignistyp, Instanz) \\
\end{longtable}
}

Diese Funktionen werden in Kapitel 6.3--6.5 direkt genutzt:
\passthrough{\lstinline!cliffs\_delta!} für den TTP-Vergleich,
\passthrough{\lstinline!cohens\_kappa!} für die
Bewertungsübereinstimmung und
\passthrough{\lstinline!aggregate\_quality\_by\_group!} für die
profilspezifische Qualitätsanalyse. Die Implementierung in einem
dedizierten \passthrough{\lstinline!utils!}-Modul (statt im Test-Code)
macht die Statistikfunktionen wiederverwendbar --- etwa für einen
zukünftigen automatisierten Scoring-Dienst im Produktivbetrieb (vgl.
Kap. 8.3.3).

Die detaillierte Auswertung aller Testmetriken, Coverage-Verteilungen
und konkreten Testbeispiele folgt in Abschnitt 6.2.

\begin{center}\rule{0.5\linewidth}{0.5pt}\end{center}

\subsection{5.7 Deployment}\label{deployment}

Das System ist als dreischichtige Cloud-Deployment auf
\textbf{Render.com} (Region: Oregon, US West) realisiert. Alle drei
Dienste werden aus demselben Repository-Branch
(\passthrough{\lstinline!deploy/render-v2!}) deployt und kommunizieren
über das private Render-Netzwerk.

\textbf{Backend (Docker Web Service)} --- Das FastAPI-Backend läuft als
Docker-Container auf einem Free-Tier-Web-Service
(\passthrough{\lstinline!0.1 vCPU!},
\passthrough{\lstinline!512 MB RAM!}). Render baut das Image aus
\passthrough{\lstinline!./backend/Dockerfile!} mit
\passthrough{\lstinline!./backend!} als Build-Kontext und startet den
Uvicorn-Server gemäß der \passthrough{\lstinline!CMD!}-Instruktion im
Dockerfile. Die Health-Check-Route \passthrough{\lstinline!GET /health!}
(vgl. Abschnitt 5.2.1) wird von Render periodisch abgefragt; bei Ausfall
wird der Dienst automatisch neu gestartet. Zugangsdaten und
Konfiguration werden als Umgebungsvariablen injiziert
(Datenbankverbindungs-URL, API-Keys für LLM-Provider). Das Backend ist
unter \passthrough{\lstinline!https://liveticker-backend.onrender.com!}
erreichbar.

\textbf{Datenbank (PostgreSQL 15)} --- Die Datenbankinstanz
\passthrough{\lstinline!liveticker-db!} läuft auf einem
Render-Managed-PostgreSQL-Dienst (Free Tier:
\passthrough{\lstinline!256 MB RAM!},
\passthrough{\lstinline!1 GB Storage!}). Die Verbindung zum Backend
erfolgt über den internen Hostnamen innerhalb des privaten
Render-Netzwerks. Beim ersten Start legt
\passthrough{\lstinline!Base.metadata.create\_all(bind=engine)!} alle
Tabellen an; nachfolgende Schema-Änderungen werden via
Alembic-Migrationen eingespielt. Da Free-Tier-Instanzen keinen
persistenten Prozess-Storage vorhalten, ist alle Zustandshaltung in der
Datenbank zentralisiert --- ein direktes Ergebnis des in Abschnitt 3.8.3
beschriebenen Stateless-Design-Prinzips.

\textbf{Frontend (Static Site)} --- Das React/TypeScript-Frontend wird
als statische Seite deployt. Render führt
\passthrough{\lstinline!cd frontend \&\& npm ci \&\& npm run build!} aus
und veröffentlicht das Verzeichnis
\passthrough{\lstinline!./frontend/build!} über sein CDN. Da kein
Server-Prozess läuft, werden die generierten Dateien direkt aus dem CDN
ausgeliefert. Das Frontend ist unter
\passthrough{\lstinline!https://liveticker-frontend.onrender.com!}
erreichbar.

\textbf{n8n (lokal mit ngrok)} --- Die n8n-Workflow-Instanz läuft lokal
und ist über einen \textbf{ngrok}-Tunnel für externe Webhooks
erreichbar. Ngrok stellt einen öffentlich adressierbaren HTTPS-Endpunkt
bereit, der eingehende Trigger-Anfragen (z. B. Spielereignisse vom
Football-API-Provider) an die lokale n8n-Instanz weiterleitet. Diese
Konfiguration ist für den Entwicklungs- und Evaluationsbetrieb
ausreichend; ein produktiver Dauerbetrieb erfordert eine persistente
n8n-Instanz mit fester Webhook-URL (z. B. self-hosted auf einem VPS oder
n8n Cloud).

Die vier Komponenten bilden zusammen die in Abschnitt 4.1 konzipierten
drei Systemschichten: n8n übernimmt die Datenbeschaffung (Schicht 1),
Backend und Datenbank die Anwendungs- und Persistenzlogik (Schicht 2),
das Frontend die Präsentation (Schicht 3). Damit ist die Implementierung
vollständig; Kapitel 6 dokumentiert die systematische Evaluation des so
realisierten Systems.
